\section{Introduction}

Many economies with fixed or heavily managed exchange rate regimes exhibit persistent gaps between the official and parallel exchange rates. These gaps reflect excess demand for foreign currency that cannot be satisfied at the official rate and instead manifests in informal or alternative markets, historically dominated by physical cash transactions. Recently, however, stablecoins---digitally transferable claims on foreign currency---have emerged as a new instrument for accessing foreign exchange in these environments.

This paper studies how the introduction of stablecoins affects the functioning of parallel foreign-exchange markets, the stability of fixed exchange rate regimes, and welfare. The central insight is that stablecoins do not merely provide a cheaper or more accessible substitute for physical dollars; they fundamentally alter the information structure of the economy. By generating high-frequency, transparent price signals, stablecoins act as a public signal of exchange-rate misalignment. This improves coordination among agents and changes the nature of speculative pressure on the regime.

We develop a model of foreign-currency demand with two margins. First, agents demand foreign currency for hedging purposes, reflecting concerns about overvaluation and future access to tradable goods. Second, agents may engage in speculative runs, acquiring foreign currency in anticipation of a regime collapse. The run decision is modeled as a global game: agents observe noisy private signals of fundamentals and a public signal whose precision depends on the scale of stablecoin usage. A crisis occurs when aggregate foreign-currency demand exceeds the authorities' capacity to defend the exchange rate.

Stablecoins affect equilibrium through two channels. The first is a cost channel: stablecoins reduce transaction costs and frictions associated with accessing foreign currency, increasing both hedge demand and speculative demand. The second is an information channel: stablecoin prices provide a precise and widely observed signal of exchange-rate misalignment, increasing the weight agents place on public information and thereby strengthening coordination in the run equilibrium.

We solve the model numerically using a fully specified global-games framework with endogenous information precision. The equilibrium jointly determines foreign-currency demand, the allocation between cash and stablecoins, the precision of the public signal, and the probability of crisis. Importantly, the model is solved without relying on reduced-form approximations: agents form beliefs using exact Bayesian updating, and the distribution of public signals is determined in a fixed point with aggregate stablecoin usage.

The model yields three main results. First, stablecoin usage increases with exchange-rate overvaluation. As fundamentals deteriorate, demand for foreign currency rises, and stablecoins---being the lower-cost instrument---capture an increasing share of this demand. Second, stablecoins amplify crisis risk through coordination. By improving the precision of public information, stablecoins increase the extent to which agents' beliefs are aligned. This strengthens the responsiveness of aggregate demand to signals and expands the region in which self-fulfilling runs occur. Third, the welfare effects of stablecoins are non-monotonic. At low levels of overvaluation, stablecoins improve welfare by reducing transaction costs and facilitating access to foreign currency. At higher levels of overvaluation, however, the amplification of coordinated runs dominates, leading to lower welfare due to more frequent and more severe crises.

This paper contributes to three strands of literature. First, it contributes to the literature on currency crises and speculative attacks by introducing an endogenous public signal whose precision is determined by financial innovation. Second, it contributes to the literature on parallel exchange rates and informal foreign-exchange markets by showing that the structure of the parallel market itself can affect equilibrium outcomes. Third, it contributes to the literature on stablecoins and digital currencies by highlighting their role not only as payment instruments but also as information-generating assets that shape macroeconomic coordination.

The remainder of the paper is organized as follows. Section 2 presents the model environment, including the information structure, foreign-currency demand, and crisis mechanism. Section 3 characterizes the equilibrium and derives the global-games cutoff structure. Section 4 discusses the role of stablecoins in shaping both costs and information precision. Section 5 presents the quantitative implementation and main results. Section 6 decomposes welfare effects across channels. Section 7 concludes with policy implications.
